\hypertarget{appendix-2fae058c82e1}{}\label{referencestart:appendix-2fae058c82e1}
\section{Exponent bias}\label{appendix-2fae058c82e1-exponent-bias}

\textbf{Curriculum depth:} Introduced · \textbf{First seen:} C004 · \textbf{Latest development:} C004

\subsection{Definition and mechanism}\label{appendix-2fae058c82e1-definition-and-mechanism}

Normal exponents are encoded as unsigned E=e+B. Binary32 B=127, so exponent 3 stores 130. Field zero and all-ones are reserved; classify before subtracting bias.

\subsection{Prerequisites}\label{appendix-2fae058c82e1-prerequisites}

\begin{itemize}
\tightlist
\item
  \hyperref[appendix-5057f9ad327f]{Binary32} --- The concrete fields provide the representation used in this mechanism.
\end{itemize}

\subsection{Used by}\label{appendix-2fae058c82e1-used-by}

\begin{itemize}
\item
  \hyperref[appendix-a1b1b4580dfd]{Floating-point addition and absorption} --- Effective exponents must be decoded before comparing scale and shifting.
\item
  \hyperref[appendix-6854d0e2218e]{Floating-point encoding and decoding} --- uses this mechanism in its explanation.
\item
  \hyperref[appendix-00a6e044c437]{Subnormals and gradual underflow} --- uses this mechanism in its explanation.
\item
  \hyperref[appendix-14bab4cab9d4]{ULP and representable spacing} --- uses this mechanism in its explanation.
\end{itemize}

\subsection{Recall and next depth}\label{appendix-2fae058c82e1-recall-and-next-depth}

Reconstruct the concrete example in the definition and identify its validity assumptions.

\textbf{Pending:} Generalize B=2\^{}(k-1)-1 to binary64 without mis-decoding subnormals.

\subsection{Source and navigation}\label{appendix-2fae058c82e1-source-and-navigation}

\hyperref[01-foundations-i-ch001-a-stretching-ruler-for-real-valued-quantities]{M004 teaching chapter} · \hyperref[appendix-fd090d982a15-m004---parts-61-76]{Per-part definitions} · \hyperref[appendix-4ad47bdd269a]{Dependency Map}
